\chapter{当前基线结果与统计证据}
\label{chap:results}

主实验结果见表~\ref{tab:protocol-a-main-results}。该表采用 resolved-audio-grouped Protocol-A，是本文推荐的主文口径，并报告每个方法在查询维度上的 bootstrap 95\% confidence interval。TRR 在 204 条查询上取得最低参数距离 8.0467，并在 Acc@0.1、Recall、Cosine 和 Module 上取得最高均值。CLAP 覆盖 201 条查询，三个查询缺少可用缓存，因此其结果只与 TRR 在共享查询上配对比较。

\reportadd{面向 TMM 修订，表~\ref{tab:protocol-a-main-results} 的角色应被严格限定为主客观证据中的 objective core table。它可以支撑“TRR 在当前已评估检索基线中更强”，但不能支撑“TRR 是最先进通用音频表示”或“系统完成端到端自动音频制作”。若主文、supplement 或 response letter 使用该表，必须同时标注 204-query split、CLAP 201-query overlap、metric definition 和 unfinished strong-baseline boundary。}

\begin{table}[H]
	\centering
		\caption{Protocol-A audio-grouped 主结果。数值为 mean [95\% CI]；Param. Dist. 越低越好，其余指标越高越好。}
		\label{tab:protocol-a-main-results}
		\scriptsize
		\resizebox{\textwidth}{!}{%
		\begin{tabular}{lrrrrrr}
			\toprule
			Method & n & Param. Dist. & Acc@0.1 & Recall & Cosine & Module \\
			\midrule
			TRR & 204 & \textbf{8.0467 [5.7805, 10.4727]} & \textbf{0.5169 [0.4860, 0.5495]} & \textbf{0.4595 [0.4232, 0.4979]} & \textbf{0.8247 [0.7820, 0.8646]} & \textbf{0.8051 [0.7681, 0.8419]} \\
			Text-RAG & 204 & 24.1705 [20.2018, 28.1073] & 0.4717 [0.4371, 0.5073] & 0.3967 [0.3592, 0.4353] & 0.4943 [0.4314, 0.5596] & 0.7456 [0.7108, 0.7809] \\
			Wav2Vec-RAG & 204 & 23.8197 [20.1623, 27.5540] & 0.4251 [0.3925, 0.4596] & 0.3505 [0.3130, 0.3893] & 0.5292 [0.4663, 0.5902] & 0.7043 [0.6668, 0.7411] \\
			FeatureNN-RAG & 204 & 19.2613 [15.9103, 22.7101] & 0.4198 [0.3896, 0.4514] & 0.3600 [0.3269, 0.3948] & 0.5479 [0.4869, 0.6089] & 0.7134 [0.6773, 0.7497] \\
			CLAP & 201 & 22.3102 [18.6405, 26.1201] & 0.4492 [0.4167, 0.4821] & 0.3971 [0.3615, 0.4329] & 0.5725 [0.5133, 0.6317] & 0.7240 [0.6884, 0.7591] \\
			\bottomrule
		\end{tabular}}
	\end{table}

	表~\ref{tab:protocol-a-significance} 汇总 TRR 相对各 baseline 的 paired significance。对于 Param. Dist.，$\Delta$ 定义为 baseline 减 TRR，正值表示 TRR 降低误差；对于其他指标，$\Delta$ 定义为 TRR 减 baseline，正值表示 TRR 提升指标。最后一列报告五个指标中最大的 Holm-corrected $p$ 值，因此它是一个保守摘要：若该值小于 0.05，则 TRR 在五个客观指标上均通过 Holm 校正后的配对检验。

	\begin{table}[H]
		\centering
		\caption{TRR 相对各 baseline 的 paired permutation significance 摘要。$\Delta$L2 为 baseline--TRR；所有检验均使用 Holm correction。}
		\label{tab:protocol-a-significance}
		\small
		\resizebox{\textwidth}{!}{%
		\begin{tabular}{lrrrl}
			\toprule
			Baseline & n & $\Delta$L2 mean [95\% CI] & $p_{\mathrm{Holm}}$(L2) & 最大 $p_{\mathrm{Holm}}$（五指标） \\
			\midrule
			Text-RAG & 204 & 16.1238 [12.3722, 19.8789] & 0.0010 & 0.0227 \\
			Wav2Vec-RAG & 204 & 15.7730 [12.2807, 19.3414] & 0.0010 & 0.0010 \\
			FeatureNN-RAG & 204 & 11.2146 [7.6011, 14.8710] & 0.0010 & 0.0011 \\
			CLAP & 201 & 14.1460 [10.3041, 18.0624] & 0.0010 & 0.0181 \\
			\bottomrule
		\end{tabular}}
	\end{table}

	表~\ref{tab:protocol-a-main-results} 与表~\ref{tab:protocol-a-significance} 共同支持一个窄而强的结论：在当前严格切分和当前已评估检索基线中，TRR 是最强的参数对齐 retrieval prior。该结论不能外推为“TRR 是通用最优音频表示”。

外推边界来自三项事实。数据集是 guitar-effects 领域数据，不覆盖全部音频任务。评价指标是参数空间 operational metrics，不能直接替代听感距离。PaSST、PANNs 和直接参数回归在当前可复核主证据中尚未完成。

Protocol-C 的结果见表~\ref{tab:protocol-c-results}。在 standard 场景，Fusion L2 为 0.3056，接近但略差于 TRR-only 的 0.2970。在 vague text 场景，Fusion 完全退回音频侧，等同 TRR-only。在 noisy audio 场景，Fusion 完全退回文本侧，避免了 TRR-only 在 embedding-space noise 下的崩溃。在 conflict 场景，Fusion 虽然优于 Text-only，但显著弱于 TRR-only。这说明当前 fusion heuristic 有可解释 fallback 行为，也有明确 failure mode。

\begin{table}[H]
	\centering
	\caption{Protocol-C 边界条件结果。权重均值来自 Fusion 分支；该协议使用 embedding-space noise，因此只作为诊断。}
	\label{tab:protocol-c-results}
	\small
	\resizebox{\textwidth}{!}{%
	\begin{tabular}{lrrrrr}
		\toprule
		\textbf{Scenario} & $\mathbb{E}[w_{\text{text}}]$ & $\mathbb{E}[w_{\text{audio}}]$ & \textbf{Text-only L2} & \textbf{TRR-only L2} & \textbf{Fusion L2} \\
		\midrule
		Standard & 0.4716 & 0.5284 & 1.5148 & \textbf{0.2970} & 0.3056 \\
		Vague text & 0.0000 & 1.0000 & 12.6517 & \textbf{0.2970} & \textbf{0.2970} \\
		Noisy audio & 1.0000 & 0.0000 & \textbf{1.5148} & 33.8549 & \textbf{1.5148} \\
		Conflict & 0.4123 & 0.5877 & 12.8572 & \textbf{0.2970} & 4.5259 \\
		\bottomrule
	\end{tabular}}
\end{table}

听测结果见表~\ref{tab:listening-summary}。Trials 2--5 中，TRR-based system 的主观评分高于 manual baseline；Trials 6--10 中，MusicGen 与 TRR-based system 的差异没有达到统计显著。由于该实验没有显式低质量 anchor，也缺少设备、参与者专业程度和人工基线时间预算等元数据，本文将其解释为 workflow relevance 的主观背景，不将其写成系统优于人工或优于生成模型的强结论。

\begin{table}[H]
	\centering
	\caption{Multiple-stimulus listening test 摘要。原始日志标签 HCAP 在论文叙述中对应 TRR-based system。}
	\label{tab:listening-summary}
	\small
	\resizebox{\textwidth}{!}{%
	\begin{tabular}{llrrrr}
		\toprule
		\textbf{Block} & \textbf{Comparison} & \textbf{n} & \textbf{Mean(A-B)} & \textbf{rbc} & \textbf{$p$ (Holm)} \\
		\midrule
		Trials 2--5 & reference vs TRR-based system & 26 & 13.78 & 0.860 & 0.00015 \\
		Trials 2--5 & TRR-based system vs manual & 26 & 19.83 & 1.000 & 0.00015 \\
		Trials 6--10 & reference vs TRR-based system & 26 & 46.71 & 0.994 & 0.00015 \\
		Trials 6--10 & MusicGen vs TRR-based system & 24 & -1.10 & 0.067 & 0.7797 \\
		\bottomrule
	\end{tabular}}
\end{table}

综合来看，当前证据支持的论文叙事应当以 Protocol-A 为中心，以 Protocol-C 说明方法边界，以听测说明用户感知背景。P0 中包含更丰富的 SOTA baseline 叙述，但缺少原始 outputs；因此它在现阶段不能替代表~\ref{tab:protocol-a-main-results} 成为主结果表。

\reportadd{TMM 论文包中，这三类结果应分工明确：Protocol-A 进入主文结果表；Protocol-C 进入边界条件或诊断段落；听测进入补充感知证据和局限讨论。任何未完成的 PaSST、PANNs、direct regression 或真实音频退化实验，只能出现在 limitations、future work 或 response 的未完成说明中，不能进入主结论。}
